\section{Lista de verificación de despliegue}
Para llevar realmente el análisis del curso al lado del servicio, los equipos suelen tener que responder varias preguntas antes de salir a producción: ¿cuál es la distribución de longitud de los prompts del negocio actual? ¿se optimiza principalmente el TTFT o el TPS? ¿el KV cache provocará fluctuaciones en la memoria de video con alta concurrencia? Si estas preguntas no se resuelven de antemano, suele ocurrir que el sistema, bajo tráfico real, muestre cuellos de botella completamente distintos a los del benchmark.

\begin{importantbox}{Antes de salir a producción, verificar al menos estas cuatro cosas}
\begin{itemize}
  \item Si las solicitudes de los usuarios consisten sobre todo en prompts largos o en respuestas largas
  \item Si el cuello de botella actual del hardware está más del lado del cómputo o del ancho de banda
  \item Si el KV cache y la estrategia de lotes corresponden a la concurrencia objetivo
  \item Si el panel de monitoreo cubre a la vez el TTFT, el TPS y el uso de memoria de video
\end{itemize}
\end{importantbox}

\subsection{Resumen de la sección}
Las diferencias teóricas entre Prefill y Decode terminan por reducirse a la lista de verificación de despliegue. Solo cuando la carga de trabajo, el hardware y las métricas se integran en un mismo proceso de verificación, la optimización de la inferencia resulta realmente ejecutable.

\newpage
